\section{Closed neutral axial dynamics and first complex mode}
\label{sec:axial-dynamics}

This section gives the current frequency-domain calculation: the radial
Einstein operator, both scalar sidebands, regular-center and radiative exterior
data, a two-sided complex-frequency match, and its remaining certification
limits.  The parity decomposition and
vacuum exterior reduce to the Regge--Wheeler framework \cite{ReggeWheeler1957},
while the coupled-channel boundary logic follows the direct matching approach
used for relativistic stars with interacting fields \cite{Macedo2016}.

\subsection{Directly linearized axial Einstein tensor}

We set $h_2=0$ and use $e^{-\ii\sigma t}$.  A direct coordinate variation of
the Einstein tensor on the general spherical metric of Eq.~\eqref{eq:metric}
gives the following coefficients of $X_A$ and $X_{AB}$ for $J=2$:
\begin{align}
\delta G_{tA}/X_A={}&-\frac{\ii\sigma}{2\gamma^2}h_1'
+\frac{\ii\sigma}{2\gamma^2}
\left(\frac{\gamma'}{\gamma}+\frac{\alpha'}{\alpha}\right)h_1
-\frac{h_0''}{2\gamma^2}
+\frac{h_0'}{2\gamma^2}
\left(\frac{\gamma'}{\gamma}+\frac{\alpha'}{\alpha}\right)
\nonumber\\
&+\frac{h_0}{\gamma^2}
\left(\frac{\alpha''}{\alpha}-
\frac{\alpha'\gamma'}{\alpha\gamma}\right)
-\frac{\ii\sigma h_1}{r\gamma^2}
-\frac{2\gamma'h_0}{r\gamma^3}
+\left(2+\frac{1}{\gamma^2}\right)\frac{h_0}{r^2},
\label{eq:direct-GtA}\\
\delta G_{rA}/X_A={}&-\frac{\sigma^2}{2\alpha^2}h_1
+\frac{\ii\sigma}{2\alpha^2}h_0'
-\frac{\ii\sigma}{r\alpha^2}h_0
\nonumber\\
&+h_1\left[
\frac{\alpha''}{\alpha\gamma^2}
-\frac{\alpha'\gamma'}{\alpha\gamma^3}
-\frac{\gamma'}{r\gamma^3}
+\frac{\alpha'}{r\alpha\gamma^2}
+\frac{2}{r^2}\right],
\label{eq:direct-GrA}\\
\delta G_{AB}/X_{AB}={}&
\frac{\ii\sigma}{\alpha^2}h_0
+\frac{1}{\gamma^2}h_1'
+\frac{1}{\gamma^2}
\left(\frac{\alpha'}{\alpha}-\frac{\gamma'}{\gamma}\right)h_1.
\label{eq:direct-GAB}
\end{align}
These expressions were obtained by differentiating the exact Christoffel and
Ricci formulas with respect to a metric bookkeeping parameter.  Their
flat-vacuum reduction is a decisive sign check.  Equations
\eqref{eq:direct-GrA} and \eqref{eq:direct-GAB} then give
\begin{align}
h_1'&=-\ii\sigma h_0,\nonumber\\
h_0'&=\frac{2h_0}{r}+
\frac{\sigma^2-4/r^2}{\ii\sigma}h_1.
\label{eq:flat-axial-pair}
\end{align}
Eliminating $h_0$ and writing $h_1=r\Psi$ yields
\begin{equation}
\Psi''+\left(\sigma^2-\frac{6}{r^2}\right)\Psi=0,
\label{eq:flat-rw-control}
\end{equation}
the $J=2$ Regge--Wheeler vacuum equation.  This test fixed the relative signs
before matter was introduced.

\subsection{Both sidebands and the complete axial stress}

At the common Fourier frequency $\sigma$, let $u_+$ multiply the perturbation
of $\Phi$ and $u_-$ the independent perturbation of $\Phi^*$.  With
$c_V=1/(2\sqrt\pi)$, the direct scalar parts of the projected stress are
\begin{align}
\tau_t={}&\frac{c_VF}{2}
\left[(2\omega+\sigma)u_--(2\omega-\sigma)u_+\right],
\label{eq:tau-t-both}\\
\tau_r={}&-\frac{\ii c_V}{2}
\left[F'(u_++u_-)-F(u_+'+u_-')\right],
\label{eq:tau-r-both}\\
\tau_2={}&\ii c_VF(u_++u_-).
\label{eq:tau-2-both}
\end{align}
The explicit metric variation in the stress tensor adds
$-\mathcal K_0h_0/2$ to the $tA$ projection and
$-\mathcal K_0h_1/2$ to the $rA$ projection, where for $\ell=1$
\begin{equation}
\mathcal K_0=\frac{1}{4\pi}\left[
-\frac{\omega^2F^2}{\alpha^2}+
\frac{F'^2}{\gamma^2}+\frac{2F^2}{r^2}+\mu^2F^2\right].
\label{eq:K0-axial}
\end{equation}
There is no such term in the $AB$ equation in Regge--Wheeler gauge because
$h_2=0$.  Equations \eqref{eq:direct-GrA}, \eqref{eq:direct-GAB}, and
\eqref{eq:tau-r-both}--\eqref{eq:tau-2-both} are two first-order Einstein
equations.  Equation \eqref{eq:direct-GtA} is retained as a dependent
constraint monitor rather than used to overdetermine the integration.

\subsection{Six-state first-order system}

Define
\begin{equation}
\bm Y=(h_0,h_1,u_+,v_+,u_-,v_-)^T,
\qquad v_\pm=u_\pm'.
\end{equation}
The tensor Einstein equation gives
\begin{equation}
h_1'=-\ii\sigma\frac{\gamma^2}{\alpha^2}h_0
+\left(\frac{\gamma'}{\gamma}-\frac{\alpha'}{\alpha}\right)h_1
+8\pi\ii c_V\gamma^2F(u_++u_-).
\label{eq:closed-h1}
\end{equation}
For compactness, introduce
\begin{equation}
\mathcal B=\frac{\alpha''}{\alpha\gamma^2}
-\frac{\alpha'\gamma'}{\alpha\gamma^3}
-\frac{\gamma'}{r\gamma^3}
+\frac{\alpha'}{r\alpha\gamma^2}+\frac{2}{r^2},
\quad
\mathcal D=F'(u_++u_-)-F(v_++v_-).
\end{equation}
The $rA$ equation becomes
\begin{equation}
h_0'=\frac{2h_0}{r}+\frac{1}{\ii\sigma}\left[
\sigma^2h_1-2\alpha^2\mathcal B h_1
-8\pi\alpha^2\mathcal K_0h_1
-8\pi\ii c_V\alpha^2\mathcal D\right].
\label{eq:closed-h0}
\end{equation}
Finally, define
\begin{equation}
\mathcal R_h=Fh_1'+\left[2F'+F
\left(\frac{\alpha'}{\alpha}-\frac{\gamma'}{\gamma}\right)\right]h_1,
\qquad \kappa_V=-\ii c_V.
\end{equation}
The two Klein--Gordon equations are
\begin{align}
v_+'={}&-\left(\frac2r+\frac{\alpha'}\alpha-\frac{\gamma'}\gamma\right)v_+
-\gamma^2\left[\frac{(\omega-\sigma)^2}{\alpha^2}-\mu^2-
\frac{6}{r^2}\right]u_+
-\frac{\ii\kappa_V\gamma^2F(2\omega-\sigma)}{\alpha^2r^2}h_0
+\frac{\kappa_V}{r^2}\mathcal R_h,
\label{eq:closed-up}\\
v_-'={}&-\left(\frac2r+\frac{\alpha'}\alpha-\frac{\gamma'}\gamma\right)v_-
-\gamma^2\left[\frac{(\omega+\sigma)^2}{\alpha^2}-\mu^2-
\frac{6}{r^2}\right]u_-
+\frac{\ii\kappa_V\gamma^2F(2\omega+\sigma)}{\alpha^2r^2}h_0
+\frac{\kappa_V}{r^2}\mathcal R_h.
\label{eq:closed-um}
\end{align}
Together with $u_\pm'=v_\pm$, Eqs.~\eqref{eq:closed-h1}--
\eqref{eq:closed-um} form the promised local system
$\bm Y'=\bm A(r,\sigma)\bm Y$.  The apparent $1/\sigma$ in
Eq.~\eqref{eq:closed-h0} reflects the chosen constraint reduction; a separate
zero-frequency formulation is required for a stationary axial perturbation.

\subsection{Center basis and channel sheets}

Let $F=f_1r+O(r^3)$ and $\alpha=\alpha_c+O(r^2)$.  The three regular
solutions have
\begin{equation}
h_0=Hr^3+O(r^5),\qquad
h_1=Br^4+O(r^6),\qquad
u_+=Pr^2+O(r^4),\qquad u_-=Nr^2+O(r^4),
\end{equation}
where
\begin{equation}
B=\frac14\left[-\frac{\ii\sigma}{\alpha_c^2}H
+8\pi\ii c_V f_1(P+N)\right].
\label{eq:axial-frobenius-B}
\end{equation}
Thus $(H,P,N)$ generate a three-column interior fundamental matrix.  At
infinity the gravitational channel is outgoing and the scalar wavenumbers are
\begin{equation}
k_\pm=\sqrt{(\omega\mp\sigma)^2-\mu^2},
\label{eq:scalar-channel-k}
\end{equation}
with the sheet chosen so that $\operatorname{Im}k_\pm\geq0$.  An open channel
is outgoing as $e^{+\ii k_\pm r_*}$; a closed channel decays on the same
sheet.  For the low-frequency mode reported below both scalar channels are
closed.

Finite-radius plane-wave conditions are inadequate for the very weakly
damped gravitational solution.  We therefore propagate $h_0/h_1$ and the two
scalar logarithmic derivatives through the exact Schwarzschild exterior.
The integration starts at $r=300$ with the exact stable logarithmic derivative
of the $J=2$ spherical Hankel function.  Linear channel pairs are evolved in
short inward segments and renormalized after every segment; this avoids both
Riccati poles and exponential overflow in a closed scalar channel.

\subsection{Two-sided matching and numerical controls}

Two-sided matching avoids the exponential contamination that affects direct
outward construction of weakly damped or massive-field modes; this is the same
boundary-value issue underlying established QNM integration and continued-
fraction methods \cite{Leaver1985,KokkotasSchmidt1999}.

Let $Y_{\rm in}$ be the $6\times3$ matrix obtained by integrating the three
regular center solutions outward and $Y_{\rm out}$ the corresponding matrix
obtained by integrating the three physical exterior channels inward.  A mode
exists when their column spaces intersect:
\begin{equation}
\det\left[Y_{\rm in}(r_m),-Y_{\rm out}(r_m)\right]=0.
\label{eq:axial-match-determinant}
\end{equation}
The smallest singular value of a Euclidean row- and column-scaled matrix is
used only as a conditioning diagnostic in the coarse search. Because those
norms contain complex conjugation, the scaled determinant is not holomorphic.
Local Newton refinement uses the unnormalized determinant in
Eq.~\eqref{eq:axial-match-determinant}. For root counting, each three-plane is
evolved in $\Lambda^3(\mathbb C^6)$: its twenty $3\times3$ minors obey the
induced compound-matrix equation, and the exterior product of the two evolved
three-forms is a holomorphic Evans function. This prevents column collapse
without introducing QR phases or Euclidean normalizations.

The production workflow searches
$0.0488\leq\operatorname{Re}\sigma\leq0.0500$ and
$-1.2\times10^{-6}\leq\operatorname{Im}\sigma\leq-10^{-7}$ on a $7\times5$
grid; no stored pole is read.  On the smaller contour
$0.04936\leq\operatorname{Re}\sigma\leq0.04944$ and
$-1.1\times10^{-6}\leq\operatorname{Im}\sigma\leq-10^{-7}$, an
eight-point-per-edge contour has winding one but a marginal maximum phase step
of $2.8622$ rad. Adaptive bisection of only under-resolved segments produces
38 boundary samples, preserves winding one, and reduces the maximum phase step
to $1.55284<\pi/2$. A stricter production rerun with a $\pi/4$ acceptance
threshold uses 46 samples and gives a maximum step $0.60974$, again with
winding one. Expanding the local contour to
$0.04930\leq\operatorname{Re}\sigma\leq0.04950$ and
$-1.5\times10^{-6}\leq\operatorname{Im}\sigma\leq10^{-7}$ gives winding one
with 43 samples and maximum phase step $0.78087<\pi/4$. These counts certify
one zero only in the declared local region; they are not a complete spectrum
or a global stability proof.

Five controls are evaluated before interpreting a root.  First, $F=0$ must
give Eq.~\eqref{eq:flat-rw-control}.  Second, the two scalar equations must
decouple to the massive flat-space $L=2$ equations.  Third, the propagated
flat exterior ratios must reproduce the exact spherical-Hankel logarithmic
derivatives.  Fourth, the root must persist when both $r_m$ and the radius at
which the stellar integration enters the Schwarzschild exterior are moved.
The first four controls pass. Fifth, the unused $tA$ Einstein equation is
evaluated on the mode using $\alpha''$ reconstructed analytically from the
equilibrium EKG equations and the chain-rule identity $Y''=A'Y+A^2Y$.
A declared fourth-order derivative-step scan lowers the relative $L_2$
defect from $2.38\times10^{-7}$ at the nominal step to
$6.85\times10^{-9}$ at multiplier 80; the smallest scanned pointwise defect
is $8.79\times10^{-8}$ at multiplier 112. The production monitor uses
Richardson extrapolation at multiplier 112. On 250, 500, and 1000 evaluation
points it gives relative $L_2=(2.82,3.21,3.56)\times10^{-8}$ and
$L_\infty=(1.55,1.67,2.19)\times10^{-7}$. The weak resolution dependence is
a documented finite-difference plateau. It controls the dependent-equation
check at the $10^{-7}$ pointwise level, although an analytic evaluation of
$A'$ would be needed to promote this numerical monitor to an exact continuum
identity.

\subsection{First frequency-domain axial mode}

For the validated $a_1^0=0.08$ background,
\begin{equation}
\omega=0.8518974044,\qquad M=1.1712711606,\qquad C_{99}=0.1048865408.
\end{equation}
Two independently refined matching setups give Table~\ref{tab:first-axial-qnm}.
\begin{table}[t]
\centering
\caption{Refined neutral $J=L=2$ axial frequency-domain pole.  The last
column is the minimum singular value of the conditioned match matrix.}
\label{tab:first-axial-qnm}
\small
\begin{tabular}{@{}ccccc@{}}
\toprule
$r_m$ & $r_e$ & $\operatorname{Re}\sigma$ & $\operatorname{Im}\sigma$
& $s_{\min}$\\
\midrule
14 & 35 & 0.0493977307853 & $-5.96519340\times10^{-7}$ & $7.91\times10^{-15}$\\
12 & 30 & 0.0493977307847 & $-5.96519314\times10^{-7}$ & $2.02\times10^{-14}$\\
\bottomrule
\end{tabular}
\end{table}
At fixed $r_{\rm far}=600$, their mean and full spread are
\begin{equation}
\sigma=0.0493977307850-5.96519327\times10^{-7}\ii,
\quad
\Delta\sigma_R=6.30\times10^{-13},
\quad
\Delta\sigma_I=2.65\times10^{-14}.
\label{eq:first-axial-qnm}
\end{equation}
The corresponding quality factor is
\begin{equation}
Q=\frac{\sigma_R}{-2\sigma_I}=4.1405\times10^4.
\end{equation}
The exterior start now uses the Schwarzschild Regge--Wheeler master expansion
and the massive-scalar Coulomb expansion through $O(r^{-3})$. For
$r_{\rm far}=120,200,300,600,900$, the imaginary parts are
$-5.9637944$, $-5.9650567$, $-5.9652382$, $-5.9651933$, and $-5.9651891$,
in units of $10^{-7}$; the 600--900 shift is $4.24\times10^{-13}$. At
$r_{\rm far}=300$, increasing the asymptotic order from two to three shifts
$(\sigma_R,\sigma_I)$ by $(3.10,4.74)\times10^{-10}$. Treating that order
shift conservatively, we quote
\begin{equation}
\sigma=0.0493977307(4)-\ii(5.9652\pm0.0048)\times10^{-7},
\qquad Q\simeq4.1405\times10^4,
\end{equation}
where the displayed uncertainty is set by asymptotic order rather than the
much smaller matching-domain spread.
The first scalar threshold is
$\sigma_{\rm th}=\mu-\omega=0.1481025957$, well above the mode.  Its scalar
tails are therefore evanescent and the negative imaginary part represents
gravitational leakage.  This is the expected signature of a long-lived,
matter-supported axial oscillation that is density-silent in the Newtonian
projection but radiatively visible in general relativity.

Equation~\eqref{eq:first-axial-qnm} is regenerated from the declared search
and establishes a counted discrete pole of the released neutral operator. The
remaining qualification is the finite-difference unused-Einstein monitor and the
absence of a genuinely independent damping calculation.

\subsection{Seven-background stable-branch continuation}

To test whether the pole is an isolated numerical feature, we continue it
through seven separately solved stable backgrounds spanning
$0.065\leq a_1^0\leq0.095$ and $0.0953\leq C_{99}\leq0.1131$. At every point,
two matching domains are refined separately and a local exterior-algebra
contour is adaptively resolved with maximum phase increment below $\pi/3$.
All seven contours have winding one, all modes remain below the first scalar
threshold, and adjacent normalized six-state profiles have overlaps above
$0.9991$. Table~\ref{tab:axial-branch} and Fig.~\ref{fig:axial-branch} therefore
establish a continuously tracked mode branch rather than one isolated root.

\begin{table}[!htbp]
\centering
\caption{Counted stable-branch continuation. The damping column is displayed
in units of $10^{-7}$; ${\cal O}_{-}$ is the overlap with the preceding
normalized six-state profile.}
\label{tab:axial-branch}
\small
\begin{tabular}{@{}rrrrrr@{}}
\toprule
$a_1^0$ & $C_{99}$ & $\operatorname{Re}\sigma$ & $-10^7\operatorname{Im}\sigma$ & $Q$ & ${\cal O}_{-}$\\
\midrule
0.065 & 0.095286 & 0.0432529917 & 3.72185 & 58107 & ---\\
0.070 & 0.098665 & 0.0453566970 & 4.41231 & 51400 & 0.999187\\
0.075 & 0.101859 & 0.0474032891 & 5.16048 & 45929 & 0.999272\\
0.080 & 0.104886 & 0.0493977307 & 5.96519 & 41405 & 0.999344\\
0.085 & 0.107760 & 0.0513442785 & 6.82511 & 37614 & 0.999406\\
0.090 & 0.110493 & 0.0532466187 & 7.73884 & 34402 & 0.999459\\
0.095 & 0.113098 & 0.0551079718 & 8.70499 & 31653 & 0.999505\\
\bottomrule
\end{tabular}
\end{table}

\begin{figure}[!htbp]
\centering
\includegraphics[width=0.94\textwidth]{figure_17_axial_branch.pdf}
\caption{Counted axial-mode branch across seven stable backgrounds. The real
frequency and gravitational damping magnitude increase smoothly with
compactness, while the quality factor decreases. Adjacent profile overlaps
near unity independently preserve mode identity.}
\label{fig:axial-branch}
\end{figure}

The largest two-domain shifts over the campaign are $4.96\times10^{-12}$ in
the real part and $8.52\times10^{-14}$ in the imaginary part. These are far
below the single-background asymptotic-order systematic. Over this interval,
an unweighted log--log fit gives the empirical relation
\begin{equation}
-\operatorname{Im}\sigma = 0.0427983\,C_{99}^{4.95770}.
\end{equation}
Its maximum and root-mean-square fractional residuals are $0.237\%$ and
$0.163\%$, respectively. This compact summary is reproducible from all seven
tabulated points, but the interval is narrow and relativistic. It is therefore
not evidence for a weak-field or post-Newtonian damping law; establishing such
a law requires continuing the branch toward the dilute limit.

\subsection{Time-domain real-frequency consistency check}

We independently rewrite the same covariant perturbation equations as a
first-order-in-time, second-order-in-space system for
\begin{equation}
(h_0,h_1,k,p,\pi,q,\chi),\qquad
k=\partial_t h_1,\quad \pi=\partial_t p,\quad \chi=\partial_t q.
\end{equation}
The $AB$ Einstein equation determines $\partial_t h_0$, the $rA$ equation
evolves $k$, and the two Klein--Gordon equations evolve $\pi$ and $\chi$.
Fourth-order finite differences, fourth-order Runge--Kutta integration, and a
sixth-derivative Kreiss--Oliger filter are used.  The stellar background is
continued by exact Schwarzschild vacuum, and the outer boundary is placed at
$r=200$, more than one gravitational wavelength from the star.  Data are
fitted only before a boundary signal can return to the extraction radius.

As an operator identity, substitution of $e^{-\ii\sigma t}$ into the
time-domain equations reproduces Eqs.~\eqref{eq:closed-h1}--
\eqref{eq:closed-um}.  With analytic radial derivatives the difference is at
machine precision.  The numerical evolution is initialized with the matched
eigenfunction but does not evaluate the shooting determinant.  Table
\ref{tab:axial-time-domain} gives the resolution and Courant study.
\begin{table}[!htbp]
\centering
\caption{Time-domain extraction at $r\simeq6$ over
$20\leq t\leq138$.  All runs use $r_{\max}=200$.}
\label{tab:axial-time-domain}
\small
\begin{tabular}{@{}ccccc@{}}
\toprule
$N_r$ & CFL & $\sigma_R^{\rm TD}$ & $\dd\ln|p|/\dd t$ & phase RMS \\
\midrule
@@TIME_DOMAIN_ROWS@@
\bottomrule
\end{tabular}
\end{table}
The regenerated time-domain matrix uses the current frequency-domain pole and
the current evolution operator at every resolution and Courant factor. The
four fitted frequencies span $4.74\times10^{-8}$; the maximum fractional error
relative to the frequency-domain value is $1.09\times10^{-6}$ and decreases to
$1.25\times10^{-7}$ on the finest grid. The phase RMS lies between
$7.90\times10^{-7}$ and $2.37\times10^{-6}$ rad. Because the initial data
are the frequency-domain eigenfunction and the run length is only
$9\times10^{-5}$ of the $1/|\sigma_I|\simeq1.68\times10^6$ damping time, the
amplitude slopes cannot measure the leakage rate or certify its sign
independently.  This experiment checks only that the time-domain operator
preserves the real oscillation frequency over the simulated window.
